\chapter{Inventaire des archives et reproduction}\label{app:inventaire}
État au 22 octobre 2026 (dernier inventaire vérifié) ; les ajouts postérieurs (E1b, identification des noyaux) sont dans le dépôt \code{mayo-a2-thesis}. Toutes les empreintes sont des SHA-256.

\section{Versions définitives à citer}
\begingroup\small\setlength{\tabcolsep}{3pt}
\begin{longtable}{@{}p{0.30\linewidth}p{0.42\linewidth}p{0.22\linewidth}@{}}
\toprule Fichier & SHA-256 & Contenu\\\midrule\endhead
\code{A2_partie1.tar.gz} & \sha{cbc0cc297ee0a0790b3cb7513d0384c6c323cfe26b6180b93fc7bd1b115477a6} & notes, code, scripts, analyses, brut S2\\
\code{A2_partie2.tar.gz} & \sha{d8dbf1b169f249c720a6e74eb0b9c6e6c9fe44eae6087ef6fe188aa21ab0fcab} & brut PH1 et CT1\\
\code{A2_partie3.tar.gz} & \sha{00d12e885de273e7851b270f8b53b780eb52fd5bbbe969b1dfe73402b687ac33} & brut S1\\
patch \code{0001} & \sha{e6c03362cbf974cf750f340453de874d667b5670a7e60eddafa54b35e29e0cea} & prototype A2\\
patch \code{0002} & \sha{751ec21fbd059e51a4b54f193f3e346a889d4b361cfdf78a3ab77c3e891b6280} & sondes\\
patch \code{0003} & \sha{b590a81d479fd4a32883a60efe4bb5525fb9f45bf92a1acd70f60b265596112d} & contrôle des sondes\\
patch \code{0004} & \sha{9c460d0de2193805f3b0c1716f62bf1390ca543778ccfb0022d995ef254a10c6} & analyse statique E1\\
{\ttfamily\scriptsize MAYO-C-A2\_source\_ e72d407.tar.gz} & \sha{e2252a0b51ab827bd097e5a0df73dcbb7ca13539af82b1e728737c58800f6833} & arbre S1/S2\\
{\ttfamily\scriptsize MAYO-C-A2-phases-v2\_ source.tar.gz} & \sha{ecb0941b086affb7189151af4b58355acb39edd4de8d7089b34ac7328cf19ef3} & arbre PH1/CT1\\
\code{E1_kit.tar.gz} & \sha{efe0b71ca5dc1553b0c0a841dc59096216a214080d988ba32731cbe2fcde18f3} & \code{run_E1.sh}, \code{e1_analyze.py}\\
\bottomrule
\end{longtable}
\endgroup
Archive E1b (sorties de compilation, dumps GCC, désassemblages) : \code{e1b_gcc16.tar.gz}, SHA-256 \sha{49452120b1b787ad2aed0f30cbcece4f7584a6063768757d92214e264dcd3faa}.\par
Archives brutes internes : S1 \code{ebc01dfe}\dots, S2 \code{b1a69769}\dots c723, PH1 \code{e1a88aec}\dots 98cc, CT1 \code{0506c2b5}\dots 92de. Les archives de la phase VM (\code{A2_donnees_et_journaux}, \code{tail_aware_etude}, \code{archive_piste_batch_EF}, etc.) sont de l'\emph{historique} et ne doivent jamais être présentées comme performances matérielles.

\section{Identification exacte du code}
Le code est déterminé par l'arbre git (indépendant de l'horodatage des commits). Base : \code{PQCMayo/MAYO-C} commit \code{74de637a7c4a7f9fafa888183bc05f5cf920cd2d}, puis \code{git am} des patchs dans l'ordre. Arbres attendus : après \code{0001} \code{87d1729c}\dots (S1, S2) ; après \code{0002} \code{c8c37c97}\dots ; après \code{0003} \code{92bcede1}\dots (PH1, CT1) ; après \code{0004} \code{bf6955da}\dots.
\begin{verbatim}
git clone https://github.com/PQCMayo/MAYO-C && cd MAYO-C
git checkout 74de637a7c4a7f9fafa888183bc05f5cf920cd2d
git am 0001-*.patch 0002-*.patch 0003-*.patch 0004-*.patch
git rev-parse HEAD^{tree}   # attendu : bf6955dae657abb087b58e13e314b946e0d283fe
\end{verbatim}

\section{Reproduire les analyses}
Prérequis : Linux, \code{python3} avec \code{numpy}, \code{tar}, \code{sha256sum}. Les trois archives dans un dossier $D$ :
\begin{verbatim}
sh reproduire_analyses.sh D travail_neuf
\end{verbatim}
Le script vérifie les empreintes, relance les neuf analyses et affiche \code{OK} si la sortie est identique (résultat obtenu : 9 sur 9). Les sorties couvrent la synthèse S1/S2 (\code{synthese_sessions.py}), la comparaison (\code{compare_sessions.py}), la décomposition (\code{phases_analyze.py}), le contraste avec/sans sondes (\code{phases_vs_ref.py}) et le contrôle CT1 (\code{controle_analyze.py}).

\section{Refaire les mesures}
Conditions : Linux, AVX2, GCC $\ge13$ (mesures : 16.2.0), cmake, \code{python3}+\code{numpy}, gouverneur \code{performance} (\code{sudo cpupower frequency-set -g performance}), cœur 1 dédié, aucun autre travail pendant l'exécution.
\begin{center}\small\begin{tabular}{@{}ll@{}}\toprule
Étape & Commande (racine de l'arbre correspondant)\\\midrule
S1, S2 & \code{SESSION=<étiquette> sh run_session_independante.sh}\\
PH1 & \code{SESSION=<étiquette> sh test/run_phases.sh}\\
CT1 & \code{SESSION=<étiquette> sh test/run_controle.sh}\\
E1 & \code{SESSION=<étiquette> sh test/run_E1.sh}\\
E1b & \code{python3 -I e1b.py <dossier_E1> <dossier_neuf>}\\\bottomrule
\end{tabular}\end{center}
Chaque \code{machine.txt} enregistre date UTC, noyau, compilateur, microcode, gouverneur, \code{boot_id} et alimentation. La reproduction des \emph{mesures} donnera des valeurs voisines et non identiques.

\chapter{Identification des noyaux dans le code généré (GCC~16.2.0, MAYO-5)}\label{app:ident}
Identification par le pas d'adresse sur $P_1$ (section~\ref{sec:protocole}). « Écritures » : écritures de 32 octets dans le corps de boucle.

\begin{table}[ht]\centering\small
\caption{Instructions / écritures de 32 octets des boucles de travail selon la construction.}
\begin{tabular}{@{}lcccc@{}}\toprule
Construction & \code{P1} off. & \code{P1} A2 & \code{P1t} off. & \code{P1t} A2\\\midrule
\code{ctx_official} & 123 / 29 & -- & 162 / 65 & --\\
\code{ctx_ab}, \code{none}, \code{esk}, \code{mcalc} & 123 / 29 & 135 / 52 & 162 / 65 & 139 / 52\\
\code{ctx_phases} (toutes sondes) & 123 / 29 & 105 / 22 & 126 / 29 & 109 / 22\\\bottomrule
\end{tabular}
\end{table}
Les différences « sans toutes les sondes » moins « toutes sondes » sont : A2 \code{P1} $+30$ (52$-$22), A2 \code{P1t} $+30$, officiel \code{P1t} $+36$ (65$-$29), officiel \code{P1} $0$.


\chapter{Exemple jouet complet de MAYO sur $\mathbb{F}_{16}$}\label{app:jouet}
Cet exemple, calculé à la main et vérifié par un programme, illustre la section~\ref{sec:schema}. Les éléments de $\mathbb{F}_{16}=\mathbb{F}_2[x]/(x^4+x+1)$ sont écrits en hexadécimal (A $=10$, B $=11$, \dots, F $=15$) ; l'addition est le OU exclusif.

\section*{Paramètres}
$n=3$, $v=2$, $o=1$, $m=2$, $k=2$ : on a bien $o<m$ et $k\,o=m$. Le secret est $O=(\mathrm{B},8)^{\mathsf T}$.

\section*{Clé publique (blocs de chaque équation $\ell=1,2$)}
\[P^{(1)}_1=\begin{pmatrix}8&9\\0&6\end{pmatrix},\ P^{(1)}_2=\begin{pmatrix}2\\2\end{pmatrix},\ P^{(1)}_3=(6);\qquad
P^{(2)}_1=\begin{pmatrix}8&\mathrm{A}\\0&8\end{pmatrix},\ P^{(2)}_2=\begin{pmatrix}\mathrm{B}\\\mathrm{C}\end{pmatrix},\ P^{(2)}_3=(\mathrm{E}).\]
Le bloc $P^{(\ell)}_3=\mathrm{Upper}\bigl(-O^{\mathsf T}(P^{(\ell)}_1O-P^{(\ell)}_2)\bigr)$ est calculé avec le secret.

\section*{Signature}
Le haché salé est supposé égal à $t=(\mathrm{F},9)$. Vinaigres tirés : $V=\begin{pmatrix}2&3\\3&\mathrm{A}\end{pmatrix}$ (lignes $v_1,v_2$). On calcule
\[Y^{(1)}=VP^{(1)}_1V^{\mathsf T}=\begin{pmatrix}8&7\\6&4\end{pmatrix},\qquad Y^{(2)}=VP^{(2)}_1V^{\mathsf T}=\begin{pmatrix}1&4\\3&9\end{pmatrix}.\]
Le second membre est $y=t+\bigl(u_{11}+z\,u_{12}+z^2u_{22}\bigr)=(8,7)$, où $z$ est la classe de $x$ modulo $z^2+z+1$ et $u_{ij}$ le $m$-vecteur des coefficients de $Y^{(\cdot)}$. La matrice du système est $A=\begin{pmatrix}2&1\\\mathrm{F}&\mathrm{D}\end{pmatrix}$ ; son déterminant vaut $6\neq0$, donc il y a une solution unique $x=(x_1,x_2)=(8,\mathrm{B})$. La signature est $S=(s_1,s_2)$ avec $s_i=(v_i+Ox_i\,\|\,x_i)$ :
\[s_1=(5,\mathrm{F},8),\qquad s_2=(\mathrm{A},\mathrm{D},\mathrm{B}).\]

\section*{Vérification}
On évalue $P^{(\ell)}(s_1)$, $P^{(\ell)}(s_2)$ et $P'^{(\ell)}(s_1,s_2)$ :
$P(s_1)=(\mathrm{A},5)$, $P(s_2)=(9,\mathrm{B})$, $P'(s_1,s_2)=(\mathrm{D},4)$. Alors
\[P^\star(S)=P(s_1)+z\,P'(s_1,s_2)+z^2P(s_2)=(\mathrm{A},5)+(6,9)+(3,5)=(\mathrm{F},9)=t.\]
La signature est acceptée. Si l'on modifie un seul élément de $s_1$, l'égalité est perdue et la signature est refusée.
